
\chapter{O processo de produção do capital e a mercadoria-energia}
\setlength{\epigraphwidth}{0.7\textwidth}

\epigraph{Havia somente um Deus, cujo nome era vapor e que falava com a voz de Malthus, McCulloch e de qualquer um que usasse máquinas.} {\textcite[p. 294]{Hobsbawm2014}}

%\epigraph{\ldots cantaremos o vibrante fervor noturno dos arsenais e dos estaleiros incendiados por violentas luas elétricas; as estações esganadas, devoradoras de serpentes que fumam; as oficinas penduradas às nuvens pelos fios contorcidos de suas fumaças; as pontes, semelhantes a ginastas gigantes que cavalgam os rios, faiscantes ao sol com um luzir de facas; os piróscafos aventurosos que farejam o horizonte, as locomotivas de largo peito, que pateiam sobre os trilhos, como enormes cavalos de aço enleados de carros; e o voo rasante dos aviões, cuja hélice freme ao vento, como uma bandeira, e parece aplaudir como uma multidão entusiasta.}{Marinetti, 1909, \emph{Manifesto Futurista}}

\epigraph{De \textit{Boulder Dam} saem algumas dezenas de cabos de cobre — longos, longos, longos cabos de cobre, talvez da espessura do seu pulso, que se estendem por centenas de quilômetros em todas as direções. Pequenos cabos de cobre transportando a potência [\textit{power}] de um rio gigantesco. E então as barras se subdividem em mais barras ... depois seguem para mais transformadores... às vezes para grandes geradores que recriam a corrente em outra forma... às vezes para motores que operam grandes maquinários industriais... para mais transformadores... depois mais subdivisões e ramificações... até que, finalmente, o rio se dispersa por toda a cidade --- acionando motores, gerando calor, produzindo luz, fazendo funcionar aparelhos diversos. O milagre de luzes quentes a partir de água fria a mais de 600 milhas de distância --- tudo isso a partir de peças de cobre e ferro dispostas de maneira específica. Motores enormes para laminar aço, ou motores minúsculos na broca de um dentista. Milhares de rodinhas girando em resposta ao movimento da grande roda em \textit{Boulder Dam}. Pare essa grande roda, e todas as rodas param; as luzes se apagam. Elas realmente estão todas conectadas.}{\textcite[Cap. 16]{Feynman2006}}


O conceito moderno de ``energia'', com que se diz que o mundo moderno ``depende de energia para funcionar'', que países possuem ``pobreza energética''\footnote{Hoje em dia, este conceito é relacionado com o 7º objetivo de desenvolvimento sustentável da ONU: ``Assegurar acesso a energia acessível, confiável, sustentável e moderna para todos''.}, etc., tem suas raízes na Revolução Industrial a partir do século XVIII, em particular no século XIX, em que os motores a vapor se tornaram o coração da ``oficina do mundo'' da época: a Inglaterra. Embora os estudos científicos à época sobre eficiência desses motores tenham se cristalizado na área da Termodinâmica --- literalmente, o estudo do ``movimento do calor'' ---, seu objetivo último sempre foi a descrição da energia enquanto entidade \emph{abstrata}, capaz de descrever cientificamente não só os motores a vapor, como também as demais fontes de força-motriz da indústria.\footnote{Não é à toa, afinal, que dois dos maiores nomes da Termodinâmica são britânicos: James Prescott Joule e Lord Kelvin (William Thomson). A relevância de ambos na Física aparece através de seus nomes no Sistema Internacional de unidades físicas: Joule é uma medida de energia, e Kelvin é uma medida de temperatura.}

Entrementes, a área do Eletromagnetismo se desenvolveu, a princípio tateando através de experimentos, e, por fim, desembocando na elegante e unificada obra de James Clerk Maxwell, cujas leis são, até hoje, a base fundamental da geração de energia elétrica: todo gerador elétrico é baseado na chamada Lei de Faraday (ou Lei de Lenz), que representa, matematicamente, como é possível obter força eletromotriz através da variação de fluxo magnético passando por um circuito fechado e de material condutor.\footnote{Isso é um exercício matemático canônico em todo manual de Física que aborde Eletromagnetismo, e.g. \textcites[pp. 312--321]{Griffiths2013}[pp. 67--71]{TongElectromag}.} A grande invenção tecnológica desta área no final do século XIX, o gerador elétrico, ainda assim está subordinada à prevalência do uso de combustíveis para gerar sua força-motriz rotacional inicial,\footnote{Este é uma variação do exercício canônico abordado na nota de rodapé acima, também presente em manuais de Física como os supracitados: se uma espira condutora gira dentro de um campo magnético uniforme, a uma velocidade angular constante, então sua força eletromotriz --- ou seja, a tensão elétrica gerada --- terá a forma de uma função periódica. Esta é a base da geração de tensão elétrica alternada, e é o alicerce de toda geração de energia, seja ela sustentável ou não --- com a exceção da energia solar, que emprega o efeito fotovoltaico (análogo, mas não idêntico, ao efeito fotoelétrico).} requerendo, portanto, mais tempo para figurar no centro do mundo moderno.

Contudo, tais teorias científicas revolucionárias não nos dizem, por si sós, a história toda. Mesmo que a intrepidez humana tenha sido, tanto antes como sempre, indispensável à descoberta e descrição adequada da natureza e de suas legalidades, ela por si só não assegurou a subsistência material desses intrépidos cientistas. Dito de outra forma: por mais que as leis naturais por eles descobertas sejam universais e anistóricas,\footnote{As legalidades \emph{naturais}, destaque-se. Aqui não se alega, como os positivistas o fazem, que as legalidades \emph{sociais} também sejam universais e anistóricas.} seus descobridores foram sujeitos singulares imersos em um contexto histórico particular: o desenvolvimento do capitalismo enquanto modo de produção mundialmente hegemônico.

O propósito do presente trabalho não é de abordar a energia enquanto conceito meramente físico, e sim enquanto um \emph{momento subordinado do capital}. A exposição acima cumpre o papel de mostrar o contexto histórico em que isso veio a ser. Cabe, agora, analisar sua importância ``sob o ponto de vista'' do capital, isto é, compreender como esta sobredeterminação do conceito ``energia'' veio a ser, como consequência necessária da dinâmica imanente do capital.

\section{Valor e capital}

O presente trabalho, portanto, tem seus pés bem fincados sobre o mundo em que a produção material busca a satisfação de necessidades \emph{de outrem}, e em que todos os produtos do trabalho chegam a seus usuários finais por meio do mercado --- em suma, a presente análise diz respeito ao mundo erigido sobre o modo de produção capitalista. Dessa forma, pressupor-se-á a extensa análise que Karl Marx fez sobre este modo de produção, em sua \emph{magnus opus}, \emph{O Capital}. Tal ``escolha teórica'', porém, não é arbitrária.

Justamente pelo ``primeiro ato histórico'' consistir na reprodução do ser humano enquanto gênero --- que pressupõe, claro, sua reprodução enquanto indivíduo singular\footnote{``\ldots{} devemos começar por constatar o primeiro pressuposto de toda a existência humana e também, portanto, de toda a história, a saber, o pressuposto de que os homens têm de estar em condições de viver para poder `fazer história'. Mas, para viver, precisa-se, antes de tudo, de comida, bebida, moradia, vestimenta e algumas coisas mais. O primeiro ato histórico é, pois, a produção dos meios para a satisfação dessas necessidades, a produção da própria vida material, e este é, sem dúvida, um ato histórico, uma condição fundamental de toda a história, que ainda hoje, assim como há milênios, tem de ser cumprida diariamente, a cada hora, simplesmente para manter os homens vivos.'' \parencite[p. 32-3]{Marx2007}} ---, é a esfera da \emph{produção} que possui prioridade ontológica na totalidade da vida humana. O ser humano devém algo além de um mero animal quando se apodera, \emph{deliberada e premeditadamente}, das legalidades naturais enquanto meios para seus fins próprios. Neste ínterim, chama-se de \emph{trabalho} o ato de se valer de tais legalidades naturais para a satisfação de suas necessidades próprias.\footnote{Isso não é uma escolha terminológica arbitrária da tradição marxista, sendo algo já reconhecido por exemplares da Economia Política clássica, por exemplo \emph{vide} John Stuart Mill: ``\ldots{} os objetos fornecidos pela natureza só servem às necessidades humanas após passarem por algum grau de transformação resultante do esforço humano. {[}\ldots{]} A extensão da transformação que as substâncias naturais sofrem antes de adquirirem a forma em que são diretamente utilizadas pelo ser humano varia desde esse grau, ou um nível ainda menor, de alteração na natureza e na aparência do objeto, até uma mudança tão completa que não resta qualquer vestígio perceptível de sua forma e estrutura originais.'' \parencite[p. 46]{Mill1909}; ``Se examinarmos qualquer outro caso do que se chama ação do homem sobre a natureza, descobriremos, da mesma forma, que as forças da natureza --- ou, em outras palavras, as propriedades da matéria --- realizam todo o trabalho, uma vez que os objetos sejam colocados na posição correta. Essa única operação --- a de colocar as coisas em locais adequados para sofrerem a ação de suas próprias forças internas e daquelas inerentes a outros objetos naturais --- é tudo o que o homem faz, ou pode fazer, com a matéria. {[}\ldots{]} O trabalho, portanto, no mundo físico, é sempre e exclusivamente empregado para colocar objetos em movimento; as propriedades da matéria e as leis da natureza encarregam-se do restante. {[}\ldots{]}'' \parencite[p. 48]{Mill1909}} É neste sentido, portanto, que Marx afirma que o trabalho é

\begin{quote}
uma condição de existência do homem, independente de todas as formas sociais, eterna necessidade natural de mediação do metabolismo entre homem e natureza \parencite[p. 120]{Marx2017}
\end{quote}

Dessa forma, o trabalho \emph{em geral} é uma característica anistórica do ser \emph{humano}:

\begin{quote}
o trabalho é antes de tudo, em termos genéticos, o ponto de partida para o tornar-se homem do homem, para a formação das suas faculdades, sendo que jamais se deve esquecer o domínio sobre si mesmo. Além do mais, o trabalho se apresenta, por um longo tempo, como o único âmbito desse desenvolvimento; todas as demais formas de atividade do homem, ligadas aos diversos valores, só se podem apresentar como autônomas depois que o trabalho atinge um nível relativamente elevado \parencite[p. 348]{Lukacs2012}
\end{quote}

A sociedade capitalista tem como particularidade o fato de o trabalho devir um \emph{momento predominante} da totalidade social, embora não o trabalho \emph{em geral}, e sim o trabalho \emph{abstrato}. Marx chama de trabalho abstrato o mero ``dispêndio de força humana de trabalho'', ``dispêndio produtivo de cérebro, músculos, nervos, mãos etc. humanos'' \parencite[p. 121]{Marx2017}, que é ``o que resta'' na comparação entre os diversos trabalhos particulares.\footnote{``Essas coisas representam apenas o fato de que em sua produção'' (das mercadorias) ``foi despendida força de trabalho humana, foi acumulado trabalho humano. Como cristais dessa substância social que lhes é comum, elas são valores -- valores de mercadorias.'' \parencite[p. 116]{Marx2017}.} Estes produtos, vistos enquanto portadores desta ``substância social comum'', possuem não somente valor de uso\footnote{``A utilidade de uma coisa faz dela um valor de uso.'' \parencite[p. 114]{Marx2017}}, como são também ``detentores'' de \emph{valor}. Quanto a isso, Marx o coloca sinteticamente:

\begin{quote}
Os homens não relacionam entre si seus produtos do trabalho como valores por considerarem essas coisas mero invólucros materiais de trabalho humano do mesmo tipo. Ao contrário. \emph{Porque} equiparam entre si seus produtos de diferentes tipos na troca, como valores, eles equiparam entre si seus diferentes trabalhos como trabalho humano. \parencite[p. 149, grifo meu]{Marx2017}
\end{quote}

Trabalho abstrato deve ser entendido, portanto, como uma \emph{abstração real}, empreendida efetivamente a todos os instantes através do metabolismo da sociedade capitalista, através das trocas de produtos dos mais variegados tipos particulares de trabalho. É na sociedade em que se \emph{produz para vender}, a ``sociedade mercantil'', que se pode entrever a substância comum da infinita variedade de trabalhos humanos, justamente porque (no limite ideal) todos estes produtos são levados ao mercado e são comparados entre si, estabelecendo proporções pelas quais pode-se trocar isso por aquilo (i.e. valores de troca). Contudo, por mais que esta abstração seja \emph{real}, efetiva na realidade concreta, sua apreensão ainda tem de ser feita \emph{idealmente}; a realidade concreta tem de ser reproduzida como um ``concreto no pensamento'' para que seja devidamente apreendida em sua totalidade. Prado sintetiza bem a apreensão da mercadoria, a ``forma elementar'' do modo de produção capitalista, enquanto uma unidade inerentemente contraditória:

\begin{quote}
O trabalho aparece de imediato como atividade específica, ou seja, como trabalho de marceneiro, pedreiro, cozinheiro etc., que produz valor{[}es{]} de uso distintos, móvel, casa, comida etc., mas não é como tal que ele comparece nos valores de troca. Como pensar o trabalho para além do trabalho concreto que produz de valor de uso? Ora, o entendimento também não pode pensar o trabalho como uma substância social{[},{]} já que para chegar a ela \emph{é preciso negar as formas aparentes para apresentar uma essência} {[}\ldots{]} Agora, a mercadoria não se apresenta mais como valor de uso e valor de troca, mas como valor de uso (aparência) e valor (essência), ou seja, ela se mostra como uma coisa inerentemente contraditória \parencite[p. 4, grifo meu]{Prado2026}.
\end{quote}

É por isso, no fim das contas, que a análise dos economistas da energia até agora abordados sempre tem de remeter à eficiência \emph{econômica} da energia: não (só) porque é a forma metodologicamente mais simples de abordar a questão da eficiência energética, mas (principalmente) porque a linguagem do mundo do capital tem como seus substantivos as mercadorias (``bens e serviços''), e como qualificativos os seus valores --- ou, em sua aparência imediata aos agentes econômicos, seus preços. O confrontamento do capital com problemas qualitativos necessariamente tem de remetê-los à sua quantificação em mercadorias e preços, justamente porque ele não pode subsumi-los senão enquanto \emph{valor}, enquanto ``massa amorfa de trabalho humano indiferenciado''.\footnote{\textcite[p. 116]{Marx2017}.}  Não pode subsumir quaisquer problemáticas sociais se elas forem antitéticas ao seu conteúdo imanente, ao seu ímpeto motriz, ao seu \emph{éthos} de autoafirmação enquanto valor que se valoriza. Assim, ao confrontar a lei do valor, a temática ecológica (assim como qualquer outra problemática) aparece como ``inadequada'' no mundo em que o capital impera como \emph{dominus}. Não é casual que este mundo seja chamado de sociedade \emph{capitalista}: seu poder e vontade dizem sempre respeito ao capital, ou melhor, \emph{são} do capital. É à \emph{ética do capital} que o mundo social passa a remeter suas atividades e mensurar sua ``adequabilidade'', ou melhor, quão ``inadequadas'' essas atividades são ou não são.\footnote{``Podemos, finalmente, concluir que, para Marx, a sociedade do capital é, na verdade, a sociedade do trabalho estranhado (capital) e de sua moral. Como em toda moral, a moral do capital envolve um conjunto de deveres ser. Como sempre, o dever ser correspondente depende da condição particular em que se encontram os indivíduos concretamente existentes. No caso da sociedade capitalista, a condição de classe, por ser determinante direto da condição econômica particular dos indivíduos, torna-se o elemento fundamental na determinação do dever ser correspondente à realização do Valor. Para a classe trabalhadora, a realização do Valor exige o comportamento adequado ao aproveitamento pelo capital (ou seja, o enfrentamento da concorrência entre os trabalhadores no mercado de trabalho). Para a classe capitalista, exige-se o comportamento adequado para a reprodução ampliada (concorrência entre capitais).'' \parencite[p. 330]{Medeiros2013a} Dessa forma, ``a ninguém, nesta forma de sociedade, é dado o direito ou a liberdade de se opor ao movimento dinâmico do capital, sob a pena da perda da condição social e, no limite, física.'' (ibid.)}

Mas o que é ``capital'', no fim das contas? No capítulo ``\emph{A transformação do dinheiro em capital}'' do livro I d'\emph{O Capital}, Marx elabora que, na circulação $D-M-D$, em que se compra ($D-M$) para vender posteriormente ($M-D$), ``mercadoria e dinheiro funcionam apenas como \emph{modos diversos de existência do próprio valor}: o dinheiro como seu modo de existência \emph{universal}, a mercadoria como seu modo de existência \emph{particular}, por assim dizer, disfarçado'' \parencite[p. 229, grifo meu]{Marx2017}; é por passar ``constantemente de uma forma a outra, sem se perder nesse movimento'' que o valor ``transforma-se em sujeito automático do processo'' (ibid., p.~229-30). Porém, quando, \emph{através deste próprio movimento}, o valor consegue devir \emph{mais-valor}, dizemos então que não é somente valor --- é, também, \emph{capital}. ``O valor se torna, assim, valor em processo, dinheiro em processo e, como tal, capital.'' \parencite[p. 231]{Marx2017}

Marx conclui que, embora o valor consiga aparecer valorizado na esfera da circulação, ele \emph{não pode} fazê-lo exclusivamente neste âmbito, pois, por mais que ele transite entre sua forma-dinheiro e sua forma-mercadoria, não é através dessas \textit{metamorfoses} --- ``mudanças de formas'' --- que ele consegue crescer em magnitude. Faz-se necessário considerar seu movimento \emph{fora desta esfera}, porém ainda dentro de seu movimento total; faz-se necessário encontrar o momento de sua dinâmica que precede sua circulação e que dela decorra como consequência: a esfera da produção. Assim, capital não é meramente dinheiro, mas sim \emph{dinheiro que busca ser mais-dinheiro}; não é meramente mercadorias, mas sim mercadorias \emph{que serão vendidas}\footnote{``O capitalista sabe que toda mercadoria, por mais miserável que seja sua aparência ou por pior que seja seu cheiro, é dinheiro, não só em sua fé, mas também na realidade.'' \parencite[p. 230]{Marx2017}}; e, no que tange à esfera da produção, não é meramente meios de produção, mas sim meios de produção \emph{para produzir mercadorias para a venda}. O capital assume a forma-dinheiro para comprar seus meios de produção;\footnote{``Como a mercadoria desaparece ao se transformar em dinheiro, neste não se percebe como ele chegou às mãos de seu possuidor ou qual mercadoria foi nele transformada. O dinheiro \emph{non olet} {[}não fede{]}, seja qual for sua origem. Se{[},{]} por um lado{[},{]} ele representa mercadoria vendida, por outro representa mercadorias compráveis.'' \parencite[p. 184]{Marx2017}} assume a forma-produção para produzir mercadorias;\footnote{``Os elementos-mercadorias {[}$FT${]} e $Mp$, que constituem o capital produtivo $P$, não possuem, como modos de existência de $P$, a mesma forma que ostentam nos diferentes mercados de mercadorias nos quais são adquiridos. Eles estão agora \emph{reunidos} e, \emph{nessa conexão, podem funcionar como capital produtivo}.'' \parencite[p. 174; grifo meu]{Marx2014}} assume a forma-mercadoria para ser vendido e retornar à forma-dinheiro, e assim sucessivamente, \emph{conquanto consiga valorizar-se mais e mais}. Em suma: o valor não pode ser capital conquanto não esteja \emph{in motu}, em expansão.

Marx chama de ``capital industrial'' o movimento do qual a circulação e a produção de mercadorias são momentos subordinados. Como Marx o coloca,

\begin{quote}
As duas formas que o valor de capital assume no interior de seus estágios de circulação são a de \emph{capital {[}dinheiro{]}} e \emph{capital-mercadoria}; sua forma própria ao estágio da produção é a de \emph{capital produtivo}. O capital, que no percurso de seu ciclo total assume e abandona de novo essas formas, cumprindo em cada uma delas sua função correspondente, é o \emph{capital industrial} --- industrial, aqui, no sentido de que ele abrange todo ramo de produção explorado de modo capitalista. \parencite[p. 131]{Marx2014}
\end{quote}

Vale destacar que capital industrial, tanto aqui quanto em Marx, \emph{não} se trata do capital \emph{no setor industrial da economia}, contraposto, por exemplo, ao capital na ``esfera financeira'', comumente chamado de ``capital financeiro''.\footnote{Inclusive, o presente trabalho prescinde completamente da controversa discussão sobre esta categoria econômica.} Trata-se da dinâmica total de produção e reprodução do capital, do processo do capital em sua totalidade, no qual as esferas da produção e da circulação de mercadorias tornam-se momentos subordinados da produção e circulação \textit{do próprio capital}. 

Porém, como asseverado no capítulo anterior, a análise do capital industrial requer que seus momentos parciais sejam estudados detidamente, a fim de que seu conteúdo possa ser apreendido em todas suas determinações e inter-relações entre seus momentos. Neste capítulo, focaremos, como Marx o fez no livro I, na esfera da produção de capital. A análise da esfera da circulação será deixada para o próximo capítulo.

% Para que o valor possa devir capital, muitas circunstâncias socioeconômicas, políticas, estruturais etc. são necessárias. \emph{Vide} Grespan, ``{[}e{]}ssencialmente, ele {[}o capital{]} é a relação social que exclui a força de trabalho da propriedade dos meios de produção e a inclui no processo de valorização como um \emph{recurso subordinado}'' \parencite[p. 49, grifo meu]{Grespan2021}. E \emph{vide} Marx:
%
% \begin{quote}
% O capital, como valor que valoriza a si mesmo, não encerra apenas relações de classes, um caráter social determinado e que repousa sobre a existência do trabalho como trabalho assalariado. Ele é um movimento, um processo cíclico que percorre diferentes estágios e, por sua vez, encerra três formas distintas do processo cíclico {[}capital-dinheiro, capital produtivo e capital-mercadoria{]}. Por isso, ele só pode ser compreendido como movimento, e não como coisa imóvel. Aqueles que consideram a autonomização do valor uma mera abstração esquecem que o movimento do capital industrial é essa mesma abstração \emph{in actu}. \parencite[p. 184]{Marx2014}
% \end{quote}
%
%

\section{O processo de produção do capital}
Para compreender o capital em seu processo de produção ``sem perturbações'', abstraímos de sua passagem pela esfera da circulação, com o que pressupomos, em particular, que tudo que ele produz seja vendido no mercado, ou seja, que seja efetivado enquanto valor. Com isso, pressupomos, para fins de análise, que a produção não se veja impedida por ``fatores exógenos''. Pressupomos, ademais, que todos seus produtos ``sejam vendidos por seus valores'', a fim de que as conclusões aqui obtidas digam respeito exclusivamente às consequências advindas \emph{deste} momento particular do capital industrial, e não, por exemplo, em decorrência a mudanças na oferta e demanda etc.

Todo capitalista sabe que a produção de mercadorias, ou ``bens e serviços'' no linguajar atual, requer a aquisição de ``fatores de produção'': máquinas, matérias-primas, contratação de força de trabalho etc. No que diz respeito à produção de capital --- i.e.~do momento de sua dinâmica na qual ele produz mercadorias e, neste fazer, valoriza seu próprio valor pressuposto ---, estes fatores tornam-se suas ``formas de existência'' neste processo. Porém, tratam-se de fatores qualitativamente diferentes no que diz respeito à produção de valor: os meios de produção que emprega --- matérias-primas, matérias auxiliares e meios de trabalho --- meramente transferem seu valor ao produto final, enquanto somente a força de trabalho que ele contrata é capaz de \emph{produzir} valor novo. Conquanto tais fatores são parte componente do processo de produção do capital, eles também agem como capital, i.e.~são partes componentes do capital-produtivo, e Marx chama-os de \emph{capital constante} e \emph{capital variável}, respectivamente.%\footnote{\textit{Vide} \parencite[p. 286]{Marx2017}: ``\ldots{} a parte do capital que se converte em meios de produção, isto é, em matérias-primas, matérias auxiliares e meios de trabalho, não altera sua grandeza de valor no processo de produção. Por essa razão, denomino-a parte constante do capital, ou, mais sucintamente: capital constante.''; ``\ldots{} a parte do capital constituída de força de trabalho modifica seu valor no processo de produção. Ela não só reproduz o equivalente de seu próprio valor, como produz um excedente, um mais-valor, que pode variar, sendo maior ou menor de acordo com as circunstâncias''.}
Como Marx o resume,
\begin{quote}
    Os mesmos componentes do capital, que, do ponto de vista do processo de trabalho, distinguem-se como fatores objetivos e subjetivos, como meios de produção e força de trabalho, distinguem-se, do ponto de vista do processo de valorização, como capital constante e capital variável. \parencite[p. 286]{Marx2017}
\end{quote}


Como a medida do valor de uma mercadoria é o tempo de trabalho socialmente necessário para a sua produção,\footnote{``Tempo de trabalho socialmente necessário é aquele requerido para produzir um valor de uso sob as \emph{condições normais} para uma dada sociedade e com o grau \emph{médio} de destreza e intensidade do trabalho.'' \parencite[p. 117, grifo meu]{Marx2017}} e como o capital constante meramente \emph{transfere} seu valor ao produto final, então é o capital variável que dita a ``variabilidade de valor'' na produção das mercadorias.

% Capital não se restringe a trabalho e meios de produção.\footnote{``\ldots não é verdade que o entendimento científico em Economia trata o capital como coisa, como máquina, instalações, etc.? E que, ao fazê-lo, concebe aquilo que se desenvolve infinitamente como algo meramente posto, limitado e finito?'' \parencite[p. 10]{Prado2009a}.}


Neste nível de abstração, não se tem em vista as particularidades destes fatores de produção, e sim que são a divisão qualitativa necessária da quantidade de valor que algum dado capital tem necessariamente de realizar a fim de empreender seu processo de reprodução. No que diz respeito à produção \textit{de valor}, o interesse do capital é de reprodução \textit{ampliada}, mas, no que diz respeito à reprodução de si próprio, enquanto capital, não há reprodução de valor que não seja ampliada, porquanto capital somente existe onde há valor que se valoriza, valor \textit{crescente} em magnitude. A importância do capital variável é, portanto, não que ele produz valor, mas sim, e principalmente, que produz \textit{mais-valor}. Como Marx o coloca assertivamente,
\begin{quote}
    Aqui [na produção de mercadorias], os valores de uso só são produzidos porque e na medida em que são o substrato material, os suportes do valor de troca. E, para nosso capitalista, trata-se de duas coisas. Primeiramente, ele quer produzir um valor de uso que tenha um valor de troca, isto é, um artigo destinado à venda, uma mercadoria. Em segundo lugar, quer produzir uma mercadoria cujo valor seja maior do que a soma do valor das mercadorias requeridas para sua produção, os meios de produção e a força de trabalho, para cuja compra ele adiantou seu dinheiro no mercado. Ele quer produzir não só um valor de uso, mas uma mercadoria; não só valor de uso, mas valor, e não só valor, mas também mais-valor. \parencite[p. 263]{Marx2017} 
\end{quote}

A produção de mercadorias em si, portanto, somente interessa ao capital porquanto são o substrato material necessário do mais-valor que ele produz para sua própria reprodução (enquanto capital). É somente neste sentido que a força de trabalho que ele emprega assume o papel, subsumida ao seu comando, de capital variável: porque ela é a forma pela qual o capital consegue produzir mais de si próprio, no que diz respeito à sua magnitude em valor. Tal magnitude, como dito acima, sempre é medida através do tempo de trabalho socialmente necessário à produção de uma dada mercadoria; todo processo de produção de mercadorias, porém, requer tanto capital constante quanto capital variável, i.e. meios de produção e força de trabalho, com o que também se deve considerar os tempos socialmente necessários a \textit{suas} produções.\footnote{``Assim, quando se considera o valor do fio [de linho], ou o tempo de trabalho requerido para sua produção, todos os diferentes processos particulares de trabalho, que, separados no tempo e no espaço, têm de ser realizados para, primeiramente, produzir o próprio algodão e a quantidade de fusos necessária à fiação e, posteriormente, para obter o fio a partir do algodão e dos fusos, \textit{podem ser considerados fases diferentes e sucessivas de um e mesmo processo de trabalho}. Todo o trabalho contido no fio é [portanto] trabalho passado.'' \parencite[p. 264-5, grifo meu]{Marx2017}} Vale relembrar que capital constante --- os meios de produção sob seu comando --- meramente \textit{transfere} seu valor ao produto final, ou seja, ele \textit{preserva} o valor do ``trabalho pretérito'' empregado no processo.\footnote{Marx se refere, \textit{en passant}, ao trabalho empregado na produção de meios de produção, comprados para um dado processo ``presente'' de produção, como tendo ocorrido ``no tempo mais-que-perfeito'' \parencite[p. 265]{Marx2017}, i.e. seus processos de produção começaram e terminaram no passado, antes do ``atual'' processo de trabalho em que serão empregados como meios de produção \textit{de facto}.} Isso é crucial, porque é pelo fato de o capital conseguir preservar seu valor inicial em seu produto que ele consegue devir um valor maior: é por preservar este valor inicial que seu ``produto de valor'' pode conter um \textit{excedente} de valor. Este excedente, é claro, tem como referência tanto o valor dos meios de produção, que é conservado no processo de produção, \textit{como também} o valor da força de trabalho, que é por ela própria \textit{re-produzido}. Como se trata de uma produção excedente de valor, Marx chama-o de \textit{mais-valor} (\textit{Mehrwert}).

Assim sendo, como a medida do valor é, em última instância, \textit{tempo}, é do interesse do capital empregar força de trabalho por quanto tempo for possível, pois ele paga pelo \textit{valor desta mercadoria}, não pelo valor \textit{advindo de seu usufruto}. O valor da força de trabalho é medido pelo valor dos meios de subsistência necessários para sua reprodução --- leia-se: a reprodução da vitalidade dos seres humanos que a vendem --- em grau \textit{normal} de qualidade, tal qual a produção de uma mercadoria qualquer.\footnote{``O valor da força de trabalho, como o de todas as outras mercadorias, é determinado pelo tempo de trabalho necessário para a produção --- e, consequentemente, também para a reprodução --- desse artigo específico. Como valor, a força de trabalho representa apenas uma quantidade determinada do trabalho social médio nela objetivado. A força de trabalho existe apenas como disposição do indivíduo vivo. A sua produção pressupõe, portanto, a existência dele. Dada a existência do indivíduo, a produção da força de trabalho consiste em sua própria reprodução ou manutenção. Para sua manutenção, o indivíduo vivo necessita de certa quantidade de meios de subsistência. Assim, o tempo de trabalho necessário à produção da força de trabalho corresponde ao tempo de trabalho necessário à produção desses meios de subsistência, ou, dito de outro modo, \textit{o valor da força de trabalho é o valor dos meios de subsistência necessários à manutenção do seu possuidor}.'' \parencite[p. 245, grifo meu]{Marx2017} Este valor, porém, não se restringe aos ``salários de subsistência'': ``As próprias necessidades naturais, como alimentação, vestimenta, aquecimento, habitação etc., são diferentes de acordo com o clima e outras peculiaridades naturais de um país. Por outro lado, a extensão das assim chamadas necessidades imediatas, assim como o modo de sua satisfação, é ela própria um produto histórico e, por isso, depende em grande medida do grau de cultura de um país, mas também, entre outros fatores, de sob quais condições e, por conseguinte, com quais costumes e exigências de vida constituiu-se a classe dos trabalhadores livres num determinado local.'' \parencite[p. 246]{Marx2017}} O trabalhador, porém, vende sua força de trabalho por um determinado período de tempo, dentro do qual o capitalista é livre para usufruí-la de acordo com seu \textit{business plan}. Isso, evidentemente, pressupõe que o valor diário da força de trabalho consiga ser reproduzido em menos tempo do que uma jornada inteira de trabalho,\footnote{Ou seja, dentro de um \textit{dia} de trabalho. ``Entre os antigos germanos, a grandeza de uma \textit{manhã} [\textit{Morgen}] de terra era medida de acordo com o trabalho de um dia e, por isso, a \textit{manhã} também era chamada de \textit{Tagwerk} [dia de trabalho] (também \textit{Tagwanne}) (\textit{jurnale} ou \textit{jurnalis}, \textit{terra jurnalis}, \textit{jornalis} ou \textit{diurnalis}) [...]'' \parencite[p. 147, n. 26]{Marx2017} Essa relação é mais imediata em francês, por exemplo, em que \textit{journée de travail} quer dizer, literalmente, \textit{dia de trabalho} --- a palavra ``dia'' se traduz tanto como \textit{jour} quanto \textit{journée}. Similarmente em italiano, \textit{giornata di lavoro} --- ``dia'' se traduz tanto como \textit{giorno} quanto \textit{giornata}, este último termo, aliás, também significando ``salário diário''.} com o que é possível dividir uma jornada de trabalho no que Marx chamou de ``tempo necessário de trabalho'' e ``tempo excedente de trabalho'':
\begin{quote}
    A soma do trabalho necessário e do mais-trabalho, isto é, dos períodos em que o trabalhador produz o valor de reposição de sua força de trabalho e o mais-valor, constitui a grandeza absoluta de seu tempo de trabalho --- a jornada de trabalho (\textit{working day}). \parencite[p. 304]{Marx2017}
\end{quote}


Diga-se de passagem: A própria jornada de trabalho, neste sentido considerado por Marx, é uma quantidade média, mesmo que venha a ser algo decretado por lei. Se a média de tempo em que o capital emprega força de trabalho, em dado setor produtivo, for, digamos, de 12 horas diárias, então uma jornada de trabalho legal de 8 horas diárias é uma mera formalidade, efetivamente letra-morta. No capítulo sobre a jornada de trabalho n'\textit{O Capital}, Marx analisa exaustivamente a conjuntura fabril da época para mostrar não só que tais limites legais eram (e são) frequentemente violados na prática efetiva dos capitalistas, como também que a jornada de trabalho ``não é, portanto, uma grandeza constante, mas variável'', cujos limites não são meramente físicos, mas também morais: 
\begin{quote}
    A variação da jornada de trabalho se move, assim, no interior de limites físicos e morais, porém ambas as formas de limites são de natureza muito elástica e permitem as mais amplas variações. Desse modo, encontramos jornadas de trabalho de 8, 10, 12, 14, 16, 18 horas, ou seja, das mais distintas durações. \parencite[p. 306]{Marx2017}
\end{quote}

Assim, tais limites da jornada de trabalho tanto são fruto da luta da classe trabalhadora, quanto são limites necessários para a reprodução \textit{em nível normal} da força de trabalho. Dessa forma, a atuação do Estado para a imposição legal de uma jornada de trabalho fixa pode ser vista tanto como uma conquista dos trabalhadores, quanto como uma restrição ``para o próprio bem'' do \textit{business as usual}.\footnote{Pode-se argumentar que a atual querela sobre a escala 6x1 no Brasil tem este caráter também. Aqui, assim como na época de Marx, tal remédio amargo tende a não ser aceito com entusiasmo pelas classes dominantes, ainda que venha a visar, em última instância, seu benefício próprio.} 

Feito este adendo, Marx aborda a questão da jornada de trabalho como sendo a categoria em que o mais-valor aparece cristalinamente como trabalho excedente (substância) e tempo excedente de trabalho (medida). É justamente a discrepância entre o valor que o capitalista paga pela compra da mercadoria força de trabalho e o valor que advém de seu usufruto o que permite a obtenção deste excedente de valor.
\begin{quote}
    A circunstância na qual a manutenção diária da força de trabalho custa [digamos] apenas meia jornada de trabalho, embora a força de trabalho possa atuar durante uma jornada inteira, e, consequentemente, o valor que ela cria durante uma jornada seja o dobro de seu próprio valor diário --- tal circunstância é, certamente, uma grande vantagem para o comprador, mas de modo algum uma injustiça para com o vendedor. \parencite[p. 270]{Marx2017}
\end{quote}

Nunca é demais destacar que este mais-valor --- assim como toda medida de valor e medidas adjacentes, como tempo necessário de trabalho --- é uma medida \textit{social}. Colocado de outra forma, tais medidas não dizem respeito à ação de capitais \textit{individuais}, e sim de um capital ``médio'', que é o referencial ao qual toda produção privada remete. Como Marx o coloca, uma jornada de trabalho é, em si, algo ``variável, não constante'', com o que o ímpeto do capital --- visto enquanto ente geral, não enquanto um dentre vários capitais singulares --- é, a princípio, aumentar o tempo diário de trabalho o quanto lhe for possível, a fim de aferir uma maior quantidade de mais-valor. Marx chama esta forma de obtenção de excedente de ``mais-valor \textit{absoluto}'': trata-se do excedente de valor produzido através do aumento da jornada de trabalho, \textit{vis-à-vis} um dado tempo necessário de trabalho; ou seja, trata-se do aumento do tempo excedente de trabalho, mantido o tempo necessário constante.\footnote{``Ao mais-valor obtido pelo prolongamento da jornada de trabalho chamo de mais-valor absoluto [...]'' \parencite[p. 390]{Marx2017}} %(Aqui já se pode entrever a ``força inexorável'' da \textit{concorrência} compelindo a ação dos capitalistas singulares; mais sobre isso adiante.)

\section{O mais-valor relativo}
Porém, além de essas mercadorias força de trabalho ``não pode[re]m ir por si mesmas ao mercado e trocar[em]-se umas pelas outras''\footnote{\textcite[p. 159]{Marx2017}.}, elas também têm inconvenientemente anexadas a si o cérebro e a dignidade humanos, o que impele o capital a ter de buscar outras formas pelas quais aferir mais-valor. O excedente de valor que é produzido em uma jornada de trabalho sempre diz respeito ao tempo trabalhado para além do tempo necessário para a reprodução da própria força de trabalho, com o que é possível obter mais-valor não somente pelo aumento extensivo dessa jornada, mas pela \textit{compressão do tempo necessário}, transformando o resto de uma dada jornada de trabalho em tempo excedente de trabalho. Este aumento ``intensivo'' de obtenção de mais-valor é chamado por Marx de ``mais-valor relativo''.\footnote{``[...] ao mais-valor que, ao contrário [do mais-valor absoluto], deriva da redução do tempo de trabalho necessário e da correspondente alteração na proporção entre as duas partes da jornada de trabalho chamo de mais-valor relativo.'' \parencite[p. 390]{Marx2017}} A obtenção de mais-valor relativo é algo que advém de mudanças no valor das mercadorias que compõem o valor da força de trabalho, o que ``retroage'' sobre o valor do capital variável que a emprega.\footnote{" ...a totalidade dos meios necessários de subsistência compõe-se de várias mercadorias, cada uma delas o produto de uma indústria distinta, e o valor de cada uma dessas mercadorias constitui uma alíquota do valor da força de trabalho. Tal valor diminui com o tempo de trabalho necessário para sua reprodução, cuja redução total é igual à soma de suas reduções em cada um dos ramos particulares da produção. Em contrapartida, nos ramos de produção que não fornecem nem meios de subsistência nem meios de produção para fabricá-los, a força produtiva aumentada deixa intocado o valor da força de trabalho." \parencite[p. 390]{Marx2017}}

Não é casual que Marx tenha deixado os famosos capítulos sobre cooperação (cap. 11), manufatura (cap. 12) e grande indústria (cap. 13) sob a seção IV de ``produção do mais-valor relativo'' do livro I d'\textit{O Capital}. Ao contrário da produção de mais-valor absoluto, que deriva do aumento da jornada de trabalho para além de um dado tempo necessário de trabalho --- i.e. da produção de valor pela força de trabalho para além da reprodução de seu próprio valor, suposto algo razoavelmente constante ---, a produção de mais-valor relativo é algo mais sutil e indiretamente aferido; enquanto a primeira pode até ser enxergada como sendo consequência da ação dos próprios capitalistas dentro de um dado setor produtivo, a última aparece aos agentes econômicos como uma consequência ``fortuita'' do vaivém do Mercado, em que se mistifica a complexa dependência dos produtos de suas empreitadas privadas. 

O capital consegue a obtenção de mais-valor relativo, segundo Marx, quando ``o prolongamento do mais-trabalho [...] resultar da \textit{redução do tempo de trabalho necessário}'', dada uma jornada de trabalho \parencite[p. 389, grifo meu]{Marx2017}. Ou seja, o excedente de valor pela via relativa é obtido não pelo aumento do tempo de trabalho, e sim pela diminuição da proporção de tempo necessário dentro de uma jornada de trabalho. A princípio, tais variações no valor da força de trabalho podem ser associados a fatores naturais: más colheitas podem aumentar o valor/preço dos alimentos, catástrofes naturais podem destruir instalações, epidemias podem reduzir a população trabalhadora etc.\footnote{No primeiro capítulo de seus \textit{Princípios de Economia Política}, Ricardo critica Smith por seu pressuposto de que o ``valor do trabalho'' era constante, justamente através do argumento de que este valor era influenciado pela variação do valor dos grãos e demais \textit{necessaries} \parencite{Ricardo2005}.} O propósito de Marx, e um de seus grandes triunfos, foi de mostrar como o próprio capital engendra e acirra legalidades sociais (tendenciais) que impulsionam tais revoluções no valor das mercadorias, legalidades estas que, ao mesmo tempo que permitem a ``superação das barreiras naturais'', tornam-se determinantes da dinâmica social. Nesse ínterim, até mesmo o famoso argumento de \textit{sunspots} feito por Jevons --- em que haveria uma relação não só de correlação, mas de \textit{causalidade} entre a ocorrência de manchas solares e as periódicas crises econômicas ao longo do século XIX \parencite{Jevons1878} --- tem de reconhecer que, por mais que fenômenos naturais possam ter uma parcela não-desprezível de causalidade quanto a fenômenos socioeconômicos, estes últimos estão inseridos na esfera \textit{social} e, dessa forma, tanto podem ser causa de fenômenos socioeconômicos ulteriores quanto podem ser acirrados por legalidades sociais vigentes. De fato, tal busca por estas causas naturais se tornam mais passíveis de investigações mais aprofundadas justamente \textit{porque} o caráter social destes fenômenos sobrepuja e encobre sua base (ineliminável) natural. 

Tais variações são, porém, algo \textit{social}, que, a princípio, não estão relacionadas com a ação dos capitalistas singulares. Dessa forma, Marx busca ilustrar como este fenômeno geral é encetado pelas ações privadas dos capitalistas individuais. Para fins de exemplo,\footnote{Cf. \textcite[pp. 391-4]{Marx2017}.} suponhamos que um dado setor tenha jornada de trabalho $J=12$h, em que $1$h de trabalho vale $\bar{f}=6$ ¢; que uma jornada de trabalho --- com a força produtiva média deste setor --- se reflita num produto de $N=12$ unidades\footnote{``Unidades'' aqui são a medida da mercadoria tratada. Marx utiliza o termo ``peça'', em linha com os exemplos usuais de fios, sapatos etc. Note-se também que, aqui, assume-se um nível mais concreto que nas discussões antecedentes: toma-se em conta a produção com mais determinações, em que o produto se subdivide em unidades individuais.}, e que cada unidade requeira o equivalente a $\bar{c}=6$ ¢ em capital constante empregado em sua produção. Dessa forma, e assumindo, como em Marx, que $12\text{ ¢} = 1\text{ sh.}$,\footnote{Em inglês, \textit{shilling}, traduzido em \textcite{Marx2017} como ``xelim''.} o valor produzido em uma jornada de trabalho é igual a
\begin{align*}
    V &= \bar{c}N + \bar{f}J = 6\text{ ¢ / un.} \cdot 12\text{un.} + 6\text{ ¢ / h} \cdot 12\text{ h} \\
    &= 144\text{ ¢}\\
    &= 12\text{ sh.}
\end{align*}

O ``valor unitário'', ou seja, o valor médio por unidade produzida, é
$$
\bar{V} \coloneqq \frac{V}{N} = 12\text{ ¢ / un.} = 1\text{ sh. / un.}
$$

\textcite{Marx2017} o denomina como ``valor individual''; prefiro o termo ``unitário'', pois remete à \textit{unidade} da mercadoria. O termo ``valor'', aqui, assume uma nova determinação: trata-se tanto do ``trabalho socialmente necessário'' que é produzido e ``se incorpora'' em valores de uso, quanto também se trata, agora, do tanto que \textit{uma unidade de mercadoria} ``contém'' deste valor total produzido; dito de outra forma, quer dizer o quanto deste valor total produzido se particiona (uniformemente) em cada unidade de mercadoria em que ele toma forma. Essa ambiguidade aparece também no mundo mais concreto: quando se fala de preços, sempre se está falando do ``preço unitário'' de algo. Quando se diz, por exemplo, que o preço de um carro é de R\$ 100.000, isso quer dizer que o preço é, na verdade, R\$ 100.000 \textit{por unidade deste tipo de carro}; para comprar, digamos, 3 carros, cada um ao mesmo preço de R\$ 100.000 por unidade, necessita-se de um equivalente monetário de ($3\text{ un.} \times$ R\$ 100.000 / un.) = R\$ 300.000. Em suma: a unidade de medida dos preços, assim como dos ``valores unitários/individuais'', é sempre de ``dinheiro por unidade''. 

Marx assume, além disso, que haja um capital específico $K$ com condições produtivas tais que ele produza, dentro de uma jornada de trabalho $J$, não $N=12$ unidades, mas $N_K = 24$ unidades.\footnote{A rigor: assumamos que a produção deste capital \textit{faça parte da média deste setor}. Ou seja, a média do setor é $N=12$ já tomando em conta que este capital específico produza $24$. Isso é o mesmo que dizer que este capital é ``excepcionalmente produtivo'', pois, para que a média seja $12$, incluindo sua produção de $24$, faz-se necessário que haja outros capitais que produzem \textit{bem abaixo} de $12$, a fim de que a média ``se estabilize'' em $12$ unidades.} Dessa forma, temos que, caso ele venda todas estas unidades --- aqui estamos abstraindo da circulação, portanto pressupondo-a ---, ele as vende por seu ``valor social'', i.e. por $1\text{ sh.}$ cada, obtendo, portanto, $24\text{ sh.}$. Como ele produziu $24$ unidades, então ele despendeu $\bar{c} \cdot N_K = 12\text{ sh.}$ em meios de produção etc., portanto tendo aferido $12\text{ sh.}$ como produto desta jornada de trabalho. Note-se que, aqui, já se pode ver o impacto deste capital mais produtivo: a média obtém $1/2\text{ sh.}/h$, ou seja, $6\text{ sh.}$ em uma jornada de 12h, ao passo que este capital obtém $1\text{ sh.}/h$, i.e. $12\text{ sh.}$ no mesmo período! Este ``excedente'' de $6 \text{ sh.}$ é o que Marx chamou de ``mais-valor adicional''.

Dado uma taxa de mais-valor $m^\prime$, ou, o que é o mesmo, a fração do tempo em que a jornada de trabalho se divide em tempo necessário --- e, portanto, em tempo excedente ---, podemos analisar como o mais-valor obtido pelo capital mais produtivo é maior que a média. Separando a jornada de trabalho em tempo necessário e tempo excedente de trabalho, da forma $J = t_N + t_E$, temos que
$$
m^\prime = \frac{t_E}{t_N} = \frac{m}{v}
$$
em que $v$ é o valor da força de trabalho e $m$ é o mais-valor aferido. Supondo que o tempo necessário seja $t_N = \nu J$, com $0 < \nu < 1$ o percentual de tempo necessário da jornada inteira, temos que
$$
m^\prime = \frac{(1-\nu) J}{\nu J} = \frac{1-\nu}{\nu}
$$

Marx pressupõe que $\nu = 10/12 = 5/6$, ou seja, $t_N = 10$h e $t_E=2$h; isso é equivalente a $m^\prime = 2/10 = 20\%$. Como o capital médio obtém $12\text{ sh.}$ ao vender $12$ unidades, dos quais $6\text{ sh.}$ são ``descontados'' como capital constante, e como cada hora de trabalho custa $6\text{ ¢ / h} = 1/2\text{ sh. / h}$, temos que, em uma jornada de $J=12$h, ele gasta $v = 5\text{ sh.}$ como capital variável, e afere $m = (6-5)=1\text{ sh.}$ como excedente de valor, i.e. mais-valor. 

O capital mais produtivo, porém, vende $24$ unidades e obtém $24\text{ sh.}$; despende $12\text{ sh.}$ em capital constante, restando-lhe $12\text{ sh.}$; seu dispêndio em capital variável é o mesmo que o da média, i.e. $5\text{ sh.}$, restando-lhe, agora, $7\text{ sh.}$ em mais-valor --- $6\text{ sh.}$ a mais que a média! Note-se, portanto, que seu mais-valor particular --- digamos, $m_K^\prime$ --- \textit{aparece} como sendo
$$
m_K^\prime = \frac{7}{5} = 140\%  \gg 20\%
$$
Equivalentemente, aparece que seu tempo necessário $t_{N, K}$ --- i.e. o tempo que este capital necessita para pagar o valor da força de trabalho que emprega --- é de
$$
\nu_K = \frac{1}{1 + m_K^\prime} = \frac{5}{12} \iff t_{N, K} = \nu_K \cdot J = 5h
$$

Ou seja, como esperado, este capital requer a \textit{metade} do tempo médio para produzir o valor equivalente da força de trabalho, justamente porque produz o \textit{dobro} da média. Essa aparência não é uma mera mistificação: este capital --- neste nível de abstração --- \textit{de fato} produz mais que a média, \textit{de fato} afere um valor maior que a média ao momento da venda\footnote{Vide nível de abstração, pressupõe-se que ele vende tudo que produz.}, e \textit{de fato} consegue recuperar o equivalente da força de trabalho que despendeu através da venda de uma proporção menor do total produzido em uma jornada de trabalho --- a média vende $5$ unidades, i.e. $5/12$ do total produzido, \textit{versus} $5/24$ deste total.\footnote{Desconsiderando as unidades vendidas com valor equivalente ao capital constante, teríamos as proporções $5/6$ e $5/12$, i.e. as proporções $\nu$ e $\nu_K$. O argumento do texto se mantém.} 


\section{Subsunção formal do trabalho pelo capital}


\section{Subsunção real do trabalho pelo capital: Maquinaria e grande indústria}
Não é uma mera coincidência que as medidas de energia na Física moderna têm termos similares aos da produção econômica \parencite{Pasquinelli2022}. \textit{Trabalho}\footnote{Em inglês, \textit{work}; em alemão, \textit{Arbeit}; em francês, \textit{travail}. Os mesmíssimos termos que aparecem na Economia são empregados (em outro contexto, evidentemente) na Física.} trata-se (\textit{grosso modo}) do emprego de uma força $F$ ao longo de alguma distância $d$, processo através do qual transfere-se energia de um meio a outro; as unidades físicas tanto do trabalho quanto da energia são as mesmas, medidas pelo Sistema Internacional (SI) em Joules (J). \textit{Potência}\footnote{Em inglês, \textit{power}; em francês, \textit{puissance} (sinônimo de ``poder''). Em alemão, \textit{Leistung}, que é também empregado no sentido de \textit{eficiência}; inclusive, é o termo que se usa para falar da ``eficiência de uma máquina'' --- não é difícil imaginar esta quantidade física historicamente sendo empregada como um \textit{proxy} da eficiência termodinâmica de motores a vapor no começo do século XIX.} trata-se desta transferência ao longo do tempo, cujas unidades são medidas, em unidades SI, em Watts (W), que é o mesmo (também pelo padrão SI) que Joules por hora. Estas unidades relacionadas a medidas de energia e de sua mudança tornaram-se quintessenciais quando da Revolução Industrial, em que a mensuração mais aderente possível dos meios de produção --- dentre os quais faziam parte, cada vez mais, combustíveis, carvão, óleo etc. --- tornava-se parte crucial da contabilidade do capitalista industrial. 

Não é ao acaso, no mesmo ínterim, que o avanço da maquinaria no Reino Unido tenha causado uma cisão entre o ``trabalho'' humano e o de uma máquina num mesmo processo de produção: o primeiro torna-se \textit{labour}, o segundo torna-se \textit{work}, ambos tendo seu ponto de contato na máquina \textit{in actu} \parencite[p. 11]{Pasquinelli2022}. Como o capitalista ``insiste em obter o que é seu'' e ``[n]ão quer ser furtado'' no que tange a seus dispêndios adiantados de capital\footnote{\textcite[p. 272]{Marx2017}.}, a mensuração mais aderente possível do processo de trabalho, em todas as facetas possíveis, torna-se um \textit{sine qua non} para sua atividade, em particular tornando-se força férrea da concorrência\footnote{A noção de ``concorrência'' em Marx tem um significado mais sutil do que sua acepção usual na Economia. Enquanto esta última tende a abordá-la como ``a quantidade de firmas em um setor de mercado'', Marx utiliza este termo para descrever a ``força social'' que --- não obstante advindo, em última instância, do fazer dos próprios capitais individuais --- acaba por impor-se a eles como ``força da natureza''. [\textbf{Descrição bem rasa; cabe ver se é cabível explicá-lo e, se sim, fazê-lo mais detidamente. Não cabe apenas numa nota de rodapé.}] É neste sentido, por exemplo, que Marx e Engels escreveram, no Manifesto Comunista, que ``Devido ao rápido aperfeiçoamento dos instrumentos de produção e ao constante progresso dos meios de comunicação, a burguesia arrasta para a torrente da civilização mesmo as nações mais bárbaras. Os baixos preços de seus produtos são a artilharia pesada que destrói todas as muralhas da China e obriga a capitularem os bárbaros mais tenazmente hostis aos estrangeiros.''.}. Dessa forma, Pasquinelli observa que

\begin{quote}
    ``as máquinas são inventadas pela organização do trabalho e, posteriormente, elas [as máquinas] acabam impondo ao trabalho suas próprias métricas. Métricas de trabalho e automação do trabalho parecem estar intimamente ligados.'' \parencite[p. 7]{Pasquinelli2022}
\end{quote}


O capitalista, por sorte, pode delegar a Física aos físicos e às conversas de bar, ocupando sua mente com o mais importante: ser um homem (ou mulher) eminentemente prático, ``que nem sempre sabe o que diz quando se encontra fora do seu negócio, mas {[}que{]} sabe muito bem o que faz dentro dele''.\footnote{\textcite[p. 269]{Marx2017}.} Os mistérios do universo não lhe importam senão como capricho; o que lhe importa, isso sim, é como tais forças naturais podem ser domadas e dobradas para dentro de seu bolso.
